\refstepcounter{section}
\section*{Lecture 31: Killing vectors for 2D metric}
\addcontentsline{toc}{section}{Lecture 31: Killing vectors for 2D metric}
In the previous section, we discussed about the Killing equation and how we can derive the \textit{isometries} of a metric from the equation. We will now see an explicit example of that, using the metric for 2D flat space,
$$\dd s^2 = \dd x^2+\dd y^2$$
From this we can find the metric and the inverse metric,
$$g_{ij} =\mqty(1 & 0 \\ 0 & 1)\qquad  g^{ij} =\mqty(1 & 0 \\ 0 & 1)$$
Intuitively, we can `see' from the metric itself, that the isometries are going to be \textit{rotations} and \textit{translations}. Let us now see whether the Killing equation says the same or not.\\[0.2cm]
Note that, since the metric is constant, the Christoffel symbols are all zero and hence covariant and partial derivatives coincide. We then have three equations:
\begin{align*}
    \partial_x \xi_x + \partial_x \xi_x &= 0\\
 \partial_y \xi_y + \partial_y \xi_y &= 0\\
  \partial_x \xi_y + \partial_y \xi_x &= 0\\
\end{align*}
The first (second) equation tells us that $\xi_x$ ($\xi_y$) is independent of $x$ ($y$) and hence we can write, 
$$\xi_x\equiv \xi_x(y)\qquad \xi_y\equiv \xi_y(x)$$
Now let us take the third equation and take the partial derivative with respect to $x$, 
$$0=\partial^2_x \xi_y + \partial_x\partial_y \xi_x = \partial^2_x \xi_y + \partial_y \cancelto{0}{\partial_x\xi_x} = \partial^2_x \xi_y \implies \boxed{\xi_y = A_2 x+B_2}$$
where $A_2$ and $B_2$ are constants to be determined. Similarly, taking partial derivative with respect to $y$ we obtain,
$$0=\partial^2_y \xi_x + \partial_y\partial_x \xi_y = \partial^2_y \xi_x + \partial_x \cancelto{0}{\partial_y\xi_y} = \partial^2_y \xi_x \implies \boxed{\xi_x = A_1 y+B_1}$$
Substituting these in the third Killing equation again gives us $A_1+A_2=0\implies A_1=-A_2$. Now, we want to characterise the vector field $\boldsymbol{\xi}$ and for that, we can write it in terms of some basis as $$\boldsymbol{\xi} = \xi^1 {\boldsymbol{\mathrm{e}}\tund{1}} +\xi^2 \boldsymbol{\mathrm{e}}\tund{2}$$
As we are using the Cartesian coordinates, a natural choice of basis is $\boldsymbol{\mathrm{e}}_i = \partial_i$ and using this, we obtain the Killing vector field as:
$$\boldsymbol{\xi} = (A_1 y+B_1) \partial_x +(-A_1x+B_2) \partial_y = B_1\partial_x + B_2\partial_y -A_1(x\partial_y-y\partial_x)$$
It is evident that the Killing vector fields are composed of three different quantities and we can now choose to focus on them one by one. 
\begin{itemize}[label=$\triangleright$]
  \item $\boxed{B_1=1, B_2=0, A_1=0}$\\[0.2cm]
  $\boldsymbol{\xi}^{(1)}=\partial_x$ is one of the infinitesimal generators of the isometry group. To see the finite transformation generated by this, we will exponentiate this. 
  $$\exp\qty(a \boldsymbol{\xi}^{(1)})=\exp\qty(a \partial_x) = \sum\limits_{n=0}^\infty \frac{a^n}{n!}\partial_x^n \equiv \mathsf{T}^1\tund{a}$$
  Now, let us act this on the coordinates:
  \begin{align*}
    \mathsf{T}^1\tund{a}(x) &= \sum\limits_{n=0}^\infty \frac{a^n}{n!}\pdv{x}{x^n}=x + a \\
\mathsf{T}^1\tund{a}(y) &= \sum\limits_{n=0}^\infty \frac{a^n}{n!}\pdv{y}{x^n}=y + 0 = y
  \end{align*}
  Thus, this generates a \textit{translation in x direction}.
  \item $\boxed{B_1=0, B_2=1, A_1=0}$\\[0.2cm]
  $\boldsymbol{\xi}^{(2)}=\partial_y$ is another infinitesimal generator of the isometry group. Similarly as before, we will exponentiate this. 
  $$\exp\qty(a \boldsymbol{\xi}^{(2)})=\exp\qty(a \partial_y) = \sum\limits_{n=0}^\infty \frac{a^n}{n!}\partial_y^n \equiv \mathsf{T}^2\tund{a}$$
  Now, let us act this on the coordinates:
  \begin{align*}
    \mathsf{T}^2\tund{a}(x) &= \sum\limits_{n=0}^\infty \frac{a^n}{n!}\pdv{x}{y^n}=x + 0 \\
\mathsf{T}^2\tund{a}(y) &= \sum\limits_{n=0}^\infty \frac{a^n}{n!}\pdv{y}{y^n}=y + a = y+a
  \end{align*}
  Thus, this generates a \textit{translation in y direction}.
   \item $\boxed{B_1=0, B_2=0, A_1=-1}$\\[0.2cm]
   $\boldsymbol{\xi}^{(3)}=x\partial_y-y\partial_x$, which is a bit non-trivial but manageable. For that, first observe that,
   $$\boldsymbol{\xi}^{(3)}(x) = -y \quad [\boldsymbol{\xi}^{(3)}]^2(x) = -x \qquad \boldsymbol{\xi}^{(3)}(y) = x\quad [\boldsymbol{\xi}^{(3)}]^2(y) = -y$$
Now, let us exponentiate the infinitesimal generator again to have some finite transformation,
\begin{align*}
  \exp\qty(\theta \boldsymbol{\xi}^{(3)}) = \sum\limits_{n=0}^\infty \frac{\theta^n}{n!}[\boldsymbol{\xi}^{(3)}]^n \equiv \mathsf{R}_\theta
\end{align*}
We can separate this into odd and even parts, since we know what happens when the generator acts odd and even number of times on a particular coordinate. 
\begin{align*}
  \mathsf{R}_\theta (x) &= \sum\limits_{n\in  \mathrm{odd}}^\infty \frac{\theta^n}{n!}[\boldsymbol{\xi}^{(3)}]^n(x) +\sum\limits_{n\in\mathrm{even}}^\infty \frac{\theta^n}{n!}[\boldsymbol{\xi}^{(3)}]^n(x)\\
  &=\sum\limits_{n=0}^\infty \frac{\theta^{2n+1}(-1)^{n+1}}{(2n+1)!}\ y+\sum\limits_{n=0}^\infty \frac{\theta^{2n}(-1)^n}{(2n)!}\ x\\
  &=x\cos\theta-y\sin\theta
\end{align*}
Similarly, acting on the coordinate y, we have:
\begin{align*}
  \mathsf{R}_\theta (y) &= \sum\limits_{n\in  \mathrm{odd}}^\infty \frac{\theta^n}{n!}[\boldsymbol{\xi}^{(3)}]^n(y) +\sum\limits_{n\in\mathrm{even}}^\infty \frac{\theta^n}{n!}[\boldsymbol{\xi}^{(3)}]^n(y)\\
  &=\sum\limits_{n=0}^\infty \frac{\theta^{2n+1}(-1)^n}{(2n+1)!}\ x+\sum\limits_{n=0}^\infty \frac{\theta^{2n}(-1)^n}{(2n)!}\ y\\
  &=x\sin\theta+y\cos\theta
\end{align*}
Thus, the finite transformation generated is rotation by some angle $\theta$. 
\end{itemize}
We see that the generators of the isometries of the metric are two translations and one rotation as was expected. The \textit{isometry group} is called the proper \textit{Euclidean group}, $\mathsf{SE}(2)$. This does not include reflections, though. If we included reflections too, the isometry group would be generalised to $\mathsf{E}(n)$, for an Euclidean space $\mathbb{E}^n$.